\documentclass[11pt]{article}
\usepackage[margin=1in]{geometry}
\usepackage{amsmath,amssymb}
\usepackage{booktabs}
\usepackage{xcolor}
\usepackage[colorlinks=true,linkcolor=black,urlcolor=blue]{hyperref}
\usepackage{parskip}
\newcommand{\xor}{\oplus}
\title{\textbf{Product-share encoding: the efficiency and hardening pass}\\[2pt]
\large What changed, and why}
\author{}
\date{2026-07-24}
\begin{document}
\maketitle

\noindent\emph{Assumes the design in} \texttt{docs/PRODUCT\_SHARE\_ENCODING}.
\emph{For the construction described from scratch, see}
\texttt{docs/NONLINEAR\_GADGETIZATION}.

\section{Summary}

The encoding worked---it was the first method to kill the computational-progress
diagonal---but it was expensive ($125\times$ the source circuit) and it carried
two weaknesses we had not looked at, both outside the reconstruction-heatmap's
field of view. This pass addresses cost and both weaknesses. The production
gadget is now \textbf{25\% smaller on 8 fewer wires with identical measured
security}, its band no longer supplies free affine invariants to a SAT attacker,
and the fold no longer advertises where each source gate begins and ends.

\begin{center}
\begin{tabular}{lll}
\toprule
 & shipped (07-23) & now (07-24) \\
\midrule
plan & $[2,2,3,3]$ & $\mathbf{[2,3,3]}$ \\
band fill & linear $\langle\alpha,x\rangle$ & \textbf{nonlinear cascade + mirror $F'$} \\
$n{=}128$ gadget & 374{,}478 gates / 576 wires & \textbf{281{,}068 / 568} \\
deg-1 ridge & depthMed 0.000, $\rho$ 0.017 & depthMed 0.000, $\rho$ 0.015 \\
deg-2 ridge & depthMed 0.000, $\rho$ $-0.035$ & depthMed 0.000, $\rho$ $-0.007$ \\
\bottomrule
\end{tabular}
\end{center}

A fourth change---a mode that removes wide gates entirely---was built and
measured, and the measurement says \textbf{do not use it}. That negative result
is \S4, and it is the most useful thing in this document.

\section{$[2,3,3]$: drop one base term}

$k$ (term count) and $\deg$ (literals per term) do different jobs. Degree hides
\textbf{algebraically}: a degree-$d$ mask sits outside the degree-$<d$
$\mathrm{GF}(2)$ span, so \emph{one} such term already forces $H=0.5$ against
every adversary of degree $<d$, whatever $k$ is. $k$ hides
\textbf{statistically}, by piling-up: $k$ stacked terms push the best affine
readout error toward $1/2$ as $1/2-2^{-(k+1)}$.

So the two degree-3 tower terms are what kill the degree-2 diagonal, and the
\emph{second} degree-2 base term buys nothing algebraic---only statistical
margin. Dropping it ($[2,2,3,3]\to[2,3,3]$) costs \textbf{24\% of the gates} and
moves the best-affine agreement from $0.57$ to $0.64$.

\textbf{The honest caveat.} That statistical cost is invisible to our
instrument: the exact-span ridge measure saturates at $H=0.5$ for any $k\ge1$,
so it scores $[2,3,3]$ and $[2,2,3,3]$ identically---as it did. Building a
statistical (best-approximation) readout is the top open item; until it exists,
every $k$-reduction is being judged by a measure blind to exactly what $k$ buys.

\section{Nonlinear band fill (and a mirror at the output)}

\textbf{The weakness.} The mask sources live on a frozen band, filled at the
input port with $\mathit{band}_j=\langle\alpha_j,x\rangle$---a \emph{linear}
function of the inputs. Every band wire was therefore a \textbf{learnable affine
invariant of the whole circuit}, held from the input port to the output port.
That is precisely the class of relation the invariant-injection experiment
showed \textbf{collapses the preimage SAT search}: forward-learned
$\mathrm{GF}(2)$ invariants turned $n{=}32$ instances from ``$>$2h timeout''
into ``$\sim$10 minutes''. We were manufacturing them by construction.

\textbf{The fix.} The fill becomes
\[
\mathit{band}_j \;=\; \mathit{junk}\;\xor\;x_{\mathrm{pivot}_j}\;\xor\;(\text{small linear part})\;\xor\;\bigoplus\nolimits^{M}(\text{two-source products})
\]
where each product's two sources are drawn from the data wires \textbf{and from
already-filled band wires}. The cascade is what makes this cheap: multiplying
two already-balanced band bits keeps the firing rate near $1/4$ while the
\emph{input degree} multiplies up the band, so high degree in $x$ is reached
without a single wide gate---every fill gate is a CNOT or a g57. (A flat
high-degree monomial would instead fire on a $2^{-d}$ fraction of inputs and
degenerate back to its linear part.)

Balance is preserved exactly, which the mask-balance argument needs: the fresh
pivot $x_{\mathrm{pivot}_j}$ is excluded from the rest of that wire's transitive
support, so the fill is a balanced function of $x$ for \emph{any} choice of the
nonlinear part. This is asserted directly in the test suite.

\textbf{Mirror $F'$.} The same fill is now emitted again after the output
bookend. The band is junk at both ports, so a two-sided composition no longer
sees it anchored only at its far end.

Cost: \textbf{free} (281{,}068 gates with the fill versus 281{,}096 without---the
difference is draw noise). The plate also improves: $\mathrm{meanH}$
$0.4931\to0.4994$, $\mathrm{stdH}$ $0.036\to0.016$.

\textbf{What this does not fix.} Band wires are still \emph{body-static}: never
written between the two ports, and heavily read. They remain trivially
identifiable statically, which is what a restriction adversary needs in order to
condition on one. The real fix is a \textbf{rolling band}---retire a band wire
mid-body, migrate its live slots with the existing re-source machinery, refill
it from current carrier values---which keeps the frozen-invariance property
piecewise. Not built.

\section{Fold-order randomization}

\textbf{The weakness.} Each source gate emitted its $(k_{\mathrm{total}}+2)^2$
fold fragments in deterministic odometer order, consecutively, all targeting the
same value's two carriers, separated from the next fold by a few RG gates. That
is a \textbf{progress clock requiring no execution at all}: an attacker segments
the fold blocks by eye and reads off the source circuit's gate boundaries---the
very thing the encoding exists to hide, handed over syntactically.

Fragments are now shuffled within each fold. This is the cheap half of the fix;
the stronger version---interleaving fragments \emph{across} adjacent folds
through the existing randomized-insertion pass, plus random ESOP re-covers and
$k$-jitter so the per-fold fragment count is not constant---is designed but not
built. Note that the final \texttt{commuting\_shuffle} and downstream mixing
already move fold material around; what was missing was randomization at the
point of emission.

\section{Narrow mode---built, measured, and \emph{not} recommended}

The motivating fear was concrete: the encoding's fold fragments reach width 6,
\texttt{--db-ctrl-cap 2} makes every gate with $\ge3$ controls permanently
ineligible for frozen-DB re-encoding, and an inert, uncrossable subpopulation
would be a fossil signature that mixing could never erase.

\texttt{--prod-max-width W} removes wide gates entirely. The mechanism is the
generalized \textbf{Barenco double sweep over dirty borrowed carriers}: a
width-$w$ conjunction becomes ${\sim}4(w-2)$ two-control gates over $w-2$
borrowed wires whose contents are \emph{arbitrary}, because each borrow is
visited an even number of times and its dirty value cancels between readings.
(The codebase already contained the $w{=}3$ instance, in
\texttt{emit\_poly\_add}'s cubic case.) So it costs \textbf{zero extra wires},
leaves no wire sitting at a constant, and preserves the gadget's ``correct under
arbitrary junk on every non-data wire'' contract.

Two implementation facts worth recording. Ladder rungs must be \textbf{exact}: a
g57's complement on a borrowed wire does \emph{not} cancel, it leaks a spurious
$c\cdot\kappa$ into the target. And an exact 2-control conjunction is
\textbf{not} a sum of g57s---every g57 carries a constant 1, so an odd number
leaves a stray 1 and an even number collapses the monomials into a plain XOR.
Rungs are therefore \texttt{comp=0} width-2 gates, which are still
DB-eligible: \texttt{db\_ctrl\_cap} filters on width, not on \texttt{comp}.

\textbf{The measurement that killed it.} A per-width commutation-leeway census
of the real $n{=}128$ gadget:

\begin{center}
\begin{tabular}{rrrr}
\toprule
width & count & mean float-box & median \\
\midrule
6 & 24{,}296 & \textbf{2778} & \textbf{2426} \\
5 & 28{,}037 & 2779 & 2415 \\
3 & 56{,}160 & 2218 & 1621 \\
4 & 65{,}907 & 1627 & 1244 \\
2 (g57) & 41{,}934 & \textbf{694} & \textbf{416} \\
\bottomrule
\end{tabular}
\end{center}

The wide mask fragments are the \textbf{most mobile gates in the circuit}---$4\times$
the mean and $5.8\times$ the median mobility of the g57s. This is the
opposite-literal separation exemption: a gate with many literals has \emph{more}
chances to hold an opposite literal against a neighbour, hence more freedom to
cross it. The premise was simply wrong. Narrowing then costs $\mathbf{6.4\times}$
the gates (1.80M vs 281k at $n{=}128$) and drops rung mobility to mean 259 /
median 96.

Security is unaffected either way---wide and narrow builds give identical ridge
results at degree 1 and degree 2, and the $[2,2]$ capability controls come back
alive ($\rho=1.00$) in \emph{both}, so the nulls are genuine hiding and the
ladder's transient partial products leak nothing.

Narrow mode's one remaining justification is \textbf{DB re-encodability}: wide
gates can never be re-encoded, which is why generation targeting reports $G=0$
forever (\S6). Whether that is worth $6.4\times$ the size is unmeasured; it needs
a \texttt{--db-dry-run --db-record} match-rate A/B on narrow versus wide
material. The flag stays in the tree, off by default.

\section{Generation targeting: what ``generation 100'' can mean here}

\texttt{G=} in the fmix report is the 5th-percentile gate generation
\textbf{over all gates}. About 62\% of a product-share gadget is width $\ge3$ and
therefore cap-ineligible: those gates are never re-encoded, their generation
stays 0, and \texttt{G=} is pinned at 0 \textbf{by construction}. ``Run to
generation $N$'' cannot mean the reported \texttt{G=} on this material, and no
amount of mixing will change that.

The achievable target is generation $N$ \textbf{over the DB-eligible material},
which is what \texttt{--gen-target} actually drives and what
\texttt{--gen-stop-frac} measures ($\mathrm{lag}/\mathrm{elig}$), with
\texttt{--gen-giveup} writing off eligible gates the DB genuinely cannot reach so
the dose stop can fire. Read \texttt{lag=/elig=} in the report; ignore
\texttt{G=}. A live run shows the split cleanly:

\begin{verbatim}
gen tgt=100 G=0 alag=280020/283384 lag=105410/106226 wlag=174598 u=12
\end{verbatim}

\texttt{wlag} is the wide population (invisible to the DB),
$\mathrm{lag}/\mathrm{elig}$ is the real progress number.

\section{Status}

Committed on \texttt{ssg-gen-mix-clean}. Library suite \textbf{171/171} (five new
tests: g57-form exactness over all polarity combinations; ladder exactness and
borrow restoration over the full input domain; share-native narrow fold; full
narrow gadget round-trip; band-fill balance and nonlinearity). Endpoint
verification passes at $n=4,16,128$.

Open, in priority order: the statistical readout (\S2); the DB match-rate A/B
(\S5); deeper fold randomization (\S4); the rolling band (\S3).

\end{document}
